% Auto-generated complete chapter from Research/deep-research-report_01_structure.md
\chapter{Auditoría por Secciones del Informe de Estrategia de Empleo}
\label{complete:research_054}

\begin{sourcemeta}
\textbf{Fuente:} \href{file:///E:/Laboral/Research/deep-research-report_01_structure.md}{Research/deep-research-report\_01\_structure.md}\\
\textbf{Seccion:} research\\
\textbf{Estado:} research\_snapshot\\
\textbf{Rol:} Research / estudio de mercado\\
\textbf{Lineas originales:} 469\\
\textbf{Ultima modificacion:} 2026-06-11 12:19\\
\textbf{Deduplicacion conservadora:} 0 bloques repetidos omitidos; 0 caracteres repetidos no reimpresos.
\end{sourcemeta}

\section*{Resumen de orientacion}
Modulo incluido completo para revision documental.

\section*{Contenido completo depurado}
\addcontentsline{toc}{section}{Contenido completo depurado}




\subsection{1. Resumen Ejecutivo}


\begin{itemize}
\item \textbf{Contenido actual:} Describe al candidato y metas salariales (USD 1.5k/3k/6k), menciona brevemente perfil (Technical Artist, multimedia, simulación) y resume las estrategias clave (investigación de mercado, posicionamiento, red de contactos, plan 30/60/90, negociación).  
\item \textbf{Faltantes / Correcciones:}  
\item \textbf{Cifras con fuente:} Incluir referencias actuales (2024–2026) para salarios y demanda (por ejemplo, Glassdoor o reportes de industria) en lugar de afirmaciones sin cita.  
\item \textbf{Prioridades claras:} Resumir al final las acciones inmediatas con datos (ej. “Optimizar CV/LinkedIn, iniciar contactos en EEUU/EU, metas de entrevistas”).  
\item \textbf{Detalle de objetivos:} Explicitar plazos y métricas (e.g. “en 3 meses >X entrevistas”).  
\item \textbf{Vincular proyectos:} Mencionar brevemente proyectos clave (TwinSight X500, ARA) como evidencia del perfil.  
\item \textbf{Lenguaje objetivo:} Evitar frases vagas o aspiracionales sin soporte (p.ej. “estrategia combinada” necesita detalle concreto).  
\end{itemize}




\subsection{2. Diagnóstico de la posición de mercado actual}


\begin{itemize}
\item \textbf{Contenido actual:} Enumera experiencia (proyectos académicos/personal), habilidades técnicas (Unity, Blender, Python, etc.), idiomas, ubicación, situación legal y metas salariales con rangos y comentarios. Incluye resumen final del perfil y metas.  
\item \textbf{Faltantes / Correcciones:}  
\item \textbf{Validar datos salariales:} Verificar con fuentes actuales salarios en EEUU/UE para Technical Artist y roles similares. Corregir rangos de USD 6.7k–10.8k mensuales según datos (Glassdoor, PayScale) para diferentes niveles.  
\item \textbf{Oferta vs demanda:} Añadir referencias de demanda de cada rol (ej. informes de empleo Unity/XR) para justificar etiquetas “Alta/Media”.  
\item \textbf{Proyectos relevantes:} Incluir que el candidato completó tesis sobre visualización técnica (TwinSight) y proyectos de IA (ARA) como experiencia práctica. Añadir su experiencia freelance (sitios web) como experiencia relevante.  
\item \textbf{Estudios incompletos:} Explicar claramente educación: “Ingeniería Multimedia (UNAD) – tesis finalizada, en espera de defensa” y “Carrera de Ingeniería Electrónica (UNAL) 80\% cursada”. Recomendar terminología (e.g. “Advanced coursework in Electronic Engineering”).  
\item \textbf{Idiomas:} Aclarar nivel exacto: Inglés C1 (certificado posible con exámenes, pruebas realizadas?), evaluar si incluir habilidades de lectura de manuales técnicos en alemán o portugués básico.  
\item \textbf{Aspectos legales y logísticos:} Precisar que como no EU, necesita contrato internacional; resumir procesos (sumario de Fe de Vida, DIAN, RUT) con referencia a fuentes oficiales.  
\item \textbf{Datos de mercado:} Añadir tendencias (p.ej. aumento de roles VR/visualización técnica en 2025) con enlaces a estudios o reportes recientes.  
\item \textbf{Competencia:} Mencionar brevemente perfil de competidores (otros junior/medios Unity tech artists) para calibrar ambición salarial.  
\item \textbf{Conclusión diagnóstica:} Reforzar que el perfil es junior pero técnicamente sólido, con gap en experiencia formal; recalcar plan de subir escalones salariales mediante proyectos y certificación.
\end{itemize}




\subsection{3. Declaración de posicionamiento recomendada}


\begin{itemize}
\item \textbf{Contenido actual:} Propuesta de enunciado principal en inglés y variantes por rol (Technical Artist Unity, 3D interactivo, XR/Digital Twin, Python/AI).  
\item \textbf{Faltantes / Correcciones:}  
\item \textbf{Enunciado clave:} Refinar el tagline único. Asegurar que incluya habilidades distintivas: \emph{“Unity Technical Artist \& Real-Time 3D Developer (Blender/URP, Shader Graph), especializado en optimización CAD-to-RT y visualización técnica interactiva con IA.”} (por ejemplo).  
\item \textbf{Coherencia de idioma:} Decidir si usar español o inglés según público objetivo. LinkedIn y CV profesional irían en inglés.  
\item \textbf{Enfatizar Blender/CAD:} Añadir Blender y optimización CAD explícitamente en el posicionamiento, dado su uso en TwinSight.  
\item \textbf{Comparación de roles:} Fundamentar la elección de roles prioritarios en evidencia: e.g. número de vacantes Unity Tech Artist vs Game TA. Incluir roles adicionales (p.ej. CAD-to-Realtime Technical Artist explícito, Tools Developer) según mercado.  
\item \textbf{Ajuste del enunciado:} Alinear con habilidades claves (C\#, Python), pero no confundir al ATS; quizá mencionar IA/ML como secundario (apoyando desarrollo).  
\item \textbf{Evaluar singulares vs combinación:} Decidir si separar “Technical Artist” y “3D Developer” en dos perfiles o unirlos. Analizar qué títulos atraen más reclutadores (usar Google Jobs o LinkedIn para frecuencia de títulos).  
\item \textbf{Coherencia con CV/LinkedIn:} Garantizar que el enunciado coincida con la declaración en resumen profesional del CV y titular de LinkedIn.  
\item \textbf{Eliminar ambigüedades:} Eliminar dobles sentidos (“destrezas en Python para producción” podría explicarse mejor, p.ej. “Automatización de pipelines con Python/IA”).
\end{itemize}




\subsection{4. Matriz de priorización de roles}


\begin{itemize}
\item \textbf{Contenido actual:} Tabla con 7 roles (Unity TA, Real-Time/WebGL Dev, XR/Simulación, Vis Tec, Game TA, Tools Dev, Python/AI Tools, Full-stack secundario) con demanda, ajuste perfil, salario, comentarios. Se concluye priorizar roles Unity/real-time, luego XR/industrial, evitar sólo games o full-stack.  
\item \textbf{Faltantes / Correcciones:}  
\item \textbf{Datos cuantitativos:} Sustentar “Demanda Alta/Media” con datos reales (p.ej. cantidad de ofertas en LinkedIn/Indeed para cada rol). Añadir referencias o al menos estimaciones de portales de empleo.  
\item \textbf{Nuevos roles:} Incluir “CAD-to-Realtime Technical Artist” como rol separado si es un nicho, y “Interactive 3D/WebGL Developer” como rol específico (actual lista mezcla con Unity TA y XR).  
\item \textbf{Remoto vs on-site:} Añadir columna sobre viabilidad remota del rol (ej. Unity TA remote es común, Game TA menos).  
\item \textbf{Freelance/contract:} Incluir columna de viabilidad contrac\_tor (p.ej. Python/AI Tools alta, dado proyectos cortos).  
\item \textbf{Localidad:} Notar qué roles demandan relocalización (XR en fabricantes, Visualización mediana).  
\item \textbf{Competencia:} Mencionar nivel de competencia: ej. Game TA alta competencia, lo que baja prioridad.  
\item \textbf{Fuente de ajuste:} Explicar cómo se calificó “ajuste al perfil” (ej. puntaje según skills).  
\item \textbf{Reordenar prioridades:} Analizar si “Digital Twin” podría ser rol separado o incluido en “Simulación Industrial”.  
\item \textbf{Detalles de salario:} Agregar rango salarial esperado por rol comparando EEUU/EU vs Colombia, con cita (como en tabla salarial).  
\item \textbf{Feedback de candidatos similares:} Ver si otros técnicos en roles similares publican experiencias o usan título específico (foros de Unity o StackOverflow Careers).  
\end{itemize}




\subsection{5. Tabla de priorización de mercados}


\begin{itemize}
\item \textbf{Contenido actual:} Tabla con regiones (EEUU/Canadá, Europa Occidental, Escandinavia, Sudeste Asiático, LATAM remota, España/Portugal, Alemania/Europa Centro, India/Middle East, Remoto global) con demanda, salario estimado y comentarios.  
\item \textbf{Faltantes / Correcciones:}  
\item \textbf{Completar filas incompletas:} La fila “Sudeste Asiático” aparece cortada (“por evid”) sin datos claros. Incluir países relevantes (India, Singapur) y rangos reales.  
\item \textbf{Agregar Australia/Nueva Zelanda:} Omitidos en la tabla; agregarlos con demanda medio-alta y salarios (Sydney/ANZ tech).  
\item \textbf{Claridad en LATAM remoto:} Diferenciar mercado local (Colombia, etc.) de empresas LATAM que contratan remotos (p.ej. Globant). Indicar salarios típicos (USD 1-3k para júnior LATAM remoto).  
\item \textbf{Salario y fuente:} Para cada región, incluir rangos salariales actualizados con referencias (reportes salary, e.g. PayScale por país). Ejemplo: USA/Can: junior Unity TA \textasciitilde{}\$70k/año (Glassdoor).  
\item \textbf{Competencia lingüística:} Notar requerimientos: e.g. en NL/Nordics idioma inglés suele bastar, en DE requeriría alemán.  
\item \textbf{Remoto global:} Específico sobre fully-remote companies (GitLab, etc. pagan por nivel global; incluir si ofrecen contratos B2B).  
\item \textbf{Formato:} Asegurar consistencia en formato (unas filas tienen “USD ≈ (junior–mid)”, otras no completan).  
\item \textbf{EOR/Contratación:} Mencionar para cada región la facilidad de contratación B2B o necesidad de EOR.  
\item \textbf{Movilidad futura:} Indicar para España/Portugal y Alemania cómo influye la potencial ciudadanía en 2 años. P.ej., España/Portugal permite residencia con EU pass.  
\item \textbf{Fuentes:} Citar rankings globales de oportunidades remotas (p.ej. reporte Stack Overflow 2025 sobre trabajos remotos).  
\item \textbf{Comentarios precisos:} Completar comentarios truncados (p.ej. añada período “mercado remoto: We Work Remotely…” a comentarios breves).  
\end{itemize}




\subsection{6. Tabla de benchmarking salarial}


\begin{itemize}
\item \textbf{Contenido actual:} Tabla por rol y mercado (EE.UU., UE, Contratista Col) con rangos en USD. Se mencionan datos de Glassdoor para Tech Artist y XR.  
\item \textbf{Faltantes / Correcciones:}  
\item \textbf{Fuentes exactas:} Para cada cifra, incluir fuente específica (por ejemplo, Glassdoor o Levels.fyi). Indicar si rango es junior/mid-level.  
\item \textbf{Distinción empleado vs contratista:} La columna “Contratista B2B Col” necesita fundamentarse en pagos freelance. Agregar comparación con salario local de ingeniería.  
\item \textbf{Salario neto vs bruto:} Especificar si USD son brutos anuales/mesuales. Aclarar que USD mensuales mencionados suponen \textasciitilde{}USD 36–144k/año según nivel.  
\item \textbf{Rangos más realistas:} Verificar si US/UE dado júnior-mid son esos rangos (parecen altos). Ajustar con datos reales (Glassdoor muestra \textasciitilde{}\$80–120k/año = \textasciitilde{}\$6.7–10k mensuales para Technical Artist senior).  
\item \textbf{Otros roles:} Incluir “Interactive 3D Developer” con WebGL (podría equipararse a “Frontend 3D Dev” en unas cifras).  
\item \textbf{Regionalización:} Podría incluir “Remoto LatAm” como columna separada.  
\item \textbf{Formato de fila:} La fila “3D WebGL/Interactive Dev” está truncada; terminar la explicación.  
\item \textbf{Contexto local:} Mencionar rango de la meta inicial (USD 1.5k) respecto a salarios locales en Colombia (p.ej. COP \$5–6M).  
\item \textbf{Actualización 2026:} Buscar publicaciones recientes sobre salarios remotos tech (ej. reportes de Deel/Remote).  
\item \textbf{Comentarios adicionales:} Explicar por qué roles de “AI/Tools” pueden superar a tech art en startups.  
\item \textbf{Presentación clara:} Tal vez dividir en dos tablas (US vs UE) para facilitar lectura.  
\end{itemize}




\subsection{7. Análisis de factibilidad de remoto desde Colombia}


\begin{itemize}
\item \textbf{Contenido actual:} Descripción del contrato B2B (facturación, impuestos, seguridad social) y opciones de pago internacional (Wise, PayPal), afiliaciones (EPS, ARL), EOR vs directo.  
\item \textbf{Faltantes / Correcciones:}  
\item \textbf{Actualización normativa:} Verificar cambios recientes en DIAN sobre facturación electrónica, retenciones y cotizaciones para independientes.  
\item \textbf{Fuentes oficiales:} Agregar enlaces a DIAN o libros tributarios para sustentar los procesos (por ejemplo, instructivo DIAN sobre exportación de servicios).  
\item \textbf{Alternativas de pago:} Incluir soluciones modernas: plataformas EOR (Deel, Remofirst), bancos internacionales (Mercantil, Bancolombia corresponsal), servicios como Payoneer, Bitwage.  
\item \textbf{Trámites clave:} Especificar pasos concretos con plazos (e.g. obtener RUT en 1 semana, afiliación ARL en X días).  
\item \textbf{Impuestos:} Detallar cómo se calcula el impuesto de renta de un contratista (tarifa máxima \textasciitilde{}39\%) y declarar facturas.  
\item \textbf{Seguridad social:} Confirmar porcentajes (EPS \textasciitilde{}12.5\%, ARL según riesgo, etc.) y plazo de pago mensual.  
\item \textbf{Afiliaciones E.U.:} Mencionar no necesaria, pero plan para seguro de salud si viaja o se muda (e.g. EHIC temporal).  
\item \textbf{Visados digitales:} Analizar si programas de visa digital nomad existen para colombianos (por ejemplo, Georgia o Portugal citizen).  
\item \textbf{Checklist:} Incluir lista de verificación: RUT, Cámara de comercio (si aplica), cuenta USD, firma electrónica para facturación.  
\item \textbf{Ejemplos / casos:} Referir estudios o guías de freelancers colombianos exportando servicios de TI.  
\item \textbf{Riesgos:} Advertir sobre fluctuación COP/USD e implicaciones en ingreso.  
\end{itemize}




\subsection{8. Estrategia futura UE (Alemania/Portugal)}


\begin{itemize}
\item \textbf{Contenido actual:} Enumera opciones: casarse con pareja alemana (permiso de residencia), obtener pasaporte portugués (ley sefardí), otras visas UE. Se menciona no asumir privilegios antes de cumplir requisitos.  
\item \textbf{Faltantes / Correcciones:}  
\item \textbf{Leyes precisas:} Citar normativa: proceso de residencia para cónyuge en Alemania (siempre y cuando casados, idioma mínimo B1); leyes portugués sefardí (tiempo de espera \textasciitilde{}2 años completando trámite) con fuente oficial.  
\item \textbf{Tiempos y requisitos:} Aclarar plazos aproximados: e.g. “en 24 meses se puede solicitar ciudadanía portuguesa si reúne documentos (certificado sefardí, antecedentes penales)”.  
\item \textbf{Visas de trabajo alternas:} Investigar visas tech de Portugal/España (p.ej. startup visa), Alemania (EU Blue Card: requiere título y X salario).  
\item \textbf{Idioma:} Incluir análisis de rentabilidad de aprender alemán vs portugués (basado en demanda de trabajo; por ej. Portugal tiene mercado IT en inglés, Alemania menos para no UE).  
\item \textbf{Mercado local DE/PT:} Mencionar oportunidades específicas (empresas que valoran portugués/EU passport).  
\item \textbf{Estrategia paso a paso:} Sugerir ruta: (1) conseguir empleo remoto estable, (2) tramitar ciudadanía portuguesa antes de emigrar, (3) casarse para visa/residencia DE si surge oportunidad.  
\item \textbf{Referencias oficiales:} Enlazar información de gobiernos sobre reunificación familiar DE y ciudadanía PT.  
\item \textbf{Riesgos:} Señalar que cambios políticos pueden afectar (ej. ley sefardí está sujeta a cambios), no confiar solo en promesas futuras.  
\item \textbf{Impacto en búsqueda:} Explicar cómo esto influye en enfocarse en roles remotos ahora sin asumir cambio inmediato de estatus legal.  
\end{itemize}




\subsection{9. Estrategia de CV}


\begin{itemize}
\item \textbf{Contenido actual:} Instrucciones para CV en inglés ATS-friendly. Se sugiere estructura (Encabezado, Resumen, Habilidades, Experiencia por proyectos, Educación, Idiomas, Cursos). Menciona brevemente proyectos TwinSight y ARA en experiencia.  
\item \textbf{Faltantes / Correcciones:}  
\item \textbf{Formato preciso:} Confirmar 1 página, con diseño claro (tipografía legible, uso moderado de negrita). Referir guías ATS para alineación de texto.  
\item \textbf{Términos en inglés:} Traducir secciones a inglés (actuales están mezcladas). Ej: “Summary” en lugar de “Resumen profesional”.  
\item \textbf{Detalles de encabezado:} Asegurar inclusión de LinkedIn/GitHub/Portfolio links activos.  
\item \textbf{Resumen profesional:} Ejemplo de párrafo concreto en inglés que resuma logros (“Completed Master's thesis on 3D drone visualization...” etc.) en 2-3 líneas.  
\item \textbf{Experiencia por proyectos:} Listar proyectos como posiciones (ej. “Independent Technical Artist”). Incluir fechas tentativas (Mar 2024–Present) para proyectos clave.  
\item \textbf{Educación:} Adaptar nomenclatura internacional: “UNAD – B.S. in Multimedia Engineering (expected 2026)”. Para UNAL: “Electronic Engineering coursework (80\% completed)”.  
\item \textbf{Cursos relevantes:} Mencionar en inglés los cursos relevantes (sin certificado) bajo sección “Professional Development” o “Courses”. E.g. “Completed CGCookie High-Fidelity Character Pipeline (self-study)”.  
\item \textbf{Palabras clave ATS:} Incluir más variantes: “Interactive 3D, Visual Computing, GPU Optimization” para cubrir más búsquedas.  
\item \textbf{Logros cuantificables:} Asegurar que cada bullet muestre impacto o métrica (porcentaje, tiempo ahorrado) como en TwinSight bullets.  
\item \textbf{Orden:} Organizar experiencias desde TwinSight (tesis) luego ARA, luego proyectos menores.  
\item \textbf{Sin jergas:} Evitar términos coloquiales; usar “3D asset optimization”, “interactive UI”, etc.  
\item \textbf{Sinónimos:} Garantizar se mencionan roles objetivo (“Unity Developer”, “Technical Artist”) repetidamente.  
\item \textbf{Error de tesis:} Asegurar que CV no describa erróneamente la tesis (debe ser TwinSight, no ARA).  
\item \textbf{Revisión bilingüe:} Aunque CV será en inglés, comprobar que el inglés sea preciso (por ejemplo, no “obtained posibility”).  
\end{itemize}




\subsection{10. Estructura completa del CV}


\begin{itemize}
\item \textbf{Análisis:} Falta detallar la estructura final con secciones y orden.  
\item \textbf{Pendientes:}  
\item Diseñar esquema definitivo de CV en inglés de 1 página con secciones: Encabezado, Perfil, Habilidades técnicas, Experiencia relevante (proyectos), Educación, Idiomas, Cursos.  
\item Definir título profesional preciso (ej. “Technical Artist \& 3D Developer”).  
\item Establecer secciones específicas: “Projects” en lugar de Experiencia laboral tradicional (pues es junior).  
\item Incluir íconos o logos (si lo permite ATS) para LinkedIn/GitHub/portafolio.  
\item Verificar formato: márgenes, uso de viñetas, consistencia.  
\item Asegurar que la línea de disponibilidad (“Full-time, willing to relocate”) esté clara.  
\item Crear ejemplo de contenido breve para cada sección (en inglés) para guiar la redacción final.  
\end{itemize}




\subsection{11. Variantes de CV (4 perfiles)}


\begin{itemize}
\item \textbf{Contenido actual:} Se menciona la necesidad de 4 variantes y su enfoque por rol (Unity TA, WebGL, XR/Digital Twin, Python/AI).  
\item \textbf{Faltantes / Correcciones:}  
\item \textbf{Definir diferencias clave:} Para cada variante, especificar qué mover hacia arriba/abajo y qué enfatizar. Ej.:  
\end{itemize}

\begin{enumerate}
\item \emph{Unity TA/3D:} poner experiencia TwinSight arriba, resaltar Unity, Shaders, Blender.  
\item \emph{WebGL/Interactive 3D:} destacar TwinSight con énfasis en WebGL, URP, UI Toolkit; incluir cualquier experiencia web (FastAPI quizás).  
\item \emph{XR/Simulación/Digital Twin:} mencionar ARA menos; resaltar prototipos industriales, XR, Bluetooth frameworks (aunque no listados aún); integrar el retrato Blender como prueba de habilidad 3D.  
\item \emph{Python/AI Tools:} Priorizar ARA Framework y AI News; posicionar proyectos de IA; mencionar brevemente habilidades web como secundaria.  
\end{enumerate}

\begin{itemize}
\item \textbf{Alineación de keywords:} Adaptar los títulos (“Technical Artist” vs “Developer”) y habilidades según el rol.  
\item \textbf{Personalización del summary:} Cambiar el resumen de perfil para cada rol destacando la meta (p.ej. “Experienced in interactive visualization…” vs “Skilled in Python ML pipelines…”).  
\item \textbf{Orden de proyectos:} Variar orden de los proyectos en Experiencia para resaltar el más relevante al rol.  
\item \textbf{Uso de plantilla:} Crear plantilla base y ajustar manualmente cada variante para mantener coherencia.  
\item \textbf{Validar ATS:} Probar cada variante con filtros simples (p.ej. pasándolo por un sitio que detecte keywords).  
\end{itemize}




\subsection{12. Banco de viñetas para CV}


\begin{itemize}
\item \textbf{Contenido actual:} Varias viñetas cuantificadas para TwinSight, ARA, FitApp y Retrato Blender.  
\item \textbf{Faltantes / Correcciones:}  
\item \textbf{Completar vacíos:} Incluir viñetas para el proyecto \emph{AI News Aggregator}. Ejemplo: “Developed full-stack AI news platform (FastAPI, React, PostgreSQL) for automated content aggregation and analysis.”  
\item \textbf{Uniformidad de idioma:} Convertir las viñetas en inglés o español consistente (preferiblemente inglés para CV internacional). Actualmente hay mezcla. Ej. “Launched personal portfolio site…” debe traducirse a inglés.  
\item \textbf{Proyectos extra:} Añadir viñetas para experiencias laborales previas (campamento YMCA, call center) si relevantes (p.ej. “Improved English proficiency by teaching activities…”).  
\item \textbf{Errores:} Corregir typos (“fitapp” vs “FitApp”, “AI News-aggregator”), asegurar consistencia (p. ej. unidades “12 participants” vs “12 participants; SUS=91.88”).  
\item \textbf{Asignación específica:} Etiquetar viñetas para cada CV variante (al pie o con notas) para fácil uso.  
\item \textbf{Lecciones aprendidas:} Incluir al menos una viñeta general que destaque la autodirección/soporte (por ej. “Self-directed learner with X hours of online training” si se desea).  
\item \textbf{Validar métricas:} Revisar que las cifras sean correctas (Triángulos, SUS scores) y citar fuente de estudio (tesis).  
\end{itemize}




\subsection{13. Estrategia completa de LinkedIn}


\begin{itemize}
\item \textbf{Contenido actual:} Sugerencias de foto, headline (varias opciones), extracto “About”, experiencia (gaps señalados), banner, actividad, keywords, boolean string ejemplo, menciones geográficas, errores actuales (tesis).  
\item \textbf{Faltantes / Correcciones:}  
\item \textbf{Titulares:} Generar las 15 opciones requeridas (todas en inglés, sin caracteres especiales) e identificar la mejor con base en claridad y keywords. Ejemplos a refinar: “Technical Artist – Unity/WebGL \& 3D Visualization” etc. Usar estudios de LinkedIn (Capgemini estudia headliners con más clics).  
\item \textbf{Headline final:} Elegir la opción que combine enfoque técnico y palabras clave. Justificar elección.  
\item \textbf{Extracto/About (en inglés):} Redactar borrador profesional de 2-3 párrafos, que equilibre perfil técnico y metas. Debe incluir: su background multimedia, pasión por visualizar datos 3D, proyectos clave y búsqueda de roles remotos. (Evitar jerga, en tono positivo/confidente.)  
\item \textbf{Experiencia actualizada:} Reescribir “Independent Technical Artist” en LinkedIn con viñetas sintéticas (en inglés) de TwinSight y ARA, utilizando bullets preparadas. Borrar errores (tesis ARA vs TwinSight). Añadir voluntariado ARCADE e idiomas.  
\item \textbf{Proyectos/Featured:} Crear secciones de “Featured” para TwinSight (con enlace GitHub/demo), ARA, Retrato. Ordenarlos según relevancia. Indicar qué añadir (videos, capturas).  
\item \textbf{Educación:} Arreglar descripción: “Thesis: TwinSight X500 interactive 3D visualization (pending defense)” en UNAD; “Electronic Eng – 80\% coursework” en UNAL. No mencionar ARA si no es tesis formal.  
\item \textbf{Orden de habilidades:} Reordenar skills destacando Unity, C\#, Python, Blender, real-time 3D al inicio.  
\item \textbf{Banner conceptual:} Diseñar idea (e.g. collage de drone wireframe + Unity interface) con texto sutil “Technical Artist – Unity \& 3D”.  
\item \textbf{Foto:} Recordar usar fondo neutro, buen encuadre, vestimenta profesional-casual (evitar selfies).  
\item \textbf{Actividad:} Planificar publicar \textasciitilde{}1/semana: reflexiones de proyecto (TwinSight), avances de cursos, artículos relevancia (Tech Art). Unirse/participar en grupos LinkedIn de Unity y XR.  
\item \textbf{Palabras clave (recruiters):} Expandir lista: “Interactive 3D Developer, CAD Visualization, Digital Twin, XR Simulation, Technical Artist Unity, Python Automation”. Incluir sinónimos en inglés.  
\item \textbf{Boolean searches:} Incluir 2–3 strings adicionales (e.g. \texttt{(“Unity” OR “Unreal”) AND “Remote” AND (colombia OR latam)}). Verificar sintaxis.  
\item \textbf{Disponibilidad:} En “About” mencionar claramente “Open to 100\% remote roles based in Colombia” y “long-term EU mobility plan”. Evitar frases ambiguas (“soon EU citizen”).  
\item \textbf{Consistencia:} Corregir o eliminar información contradictoria (p.ej. portada de LinkedIn actual, proyectos en educación).  
\item \textbf{Idioma:} Decidir el idioma principal (en inglés para audiencia internacional). Si el perfil está en inglés, toda la info (actividades, comentarios) en inglés.  
\end{itemize}




\subsection{14. Auditoría y plan de acción de GitHub}


\begin{itemize}
\item \textbf{Contenido actual:} Recomendaciones para repos clave (TwinSight, ARA, ai-news, blender portrait), públicos vs privados, README recruiter-friendly, GitHub Pages, etiquetado, actividad, limpieza de código, manejo de AI.  
\item \textbf{Faltantes / Correcciones:}  
\item \textbf{Repositorios faltantes:} Revisar repos mencionados en plan pero no listados: \emph{deus-ex-machina}, \emph{skills-introduction-to-github}, otros en GitHub (buscar privados pendientes). Decidir status (publicar/archivar).  
\item \textbf{Priorizar destacados:} Confirmar 3–5 repos con mayor impacto: p.ej. TwinSight (portfolio principal), ARA Framework, Blender portrait (si contiene procesos), \emph{ai-news-aggregator} (si es mínimo viable), \emph{3Dweb}.  
\item \textbf{Readmes:} Escribir README detallado para TwinSight y ARA con enfoque en uso (no solo código). Incluir capturas/gifs si es posible y enlace a demo WebGL. Para repos con código incompleto, aclarar “proyecto en desarrollo”.  
\item \textbf{GitHub Pages:} Generar un sitio estático (delarge95.github.io) que sirva de portfolio mínimo, enlazando proyectos.  
\item \textbf{Estandarizar etiquetas:} Usar tags (archivos \texttt{.github/release.toml} o en descripción) que incluyan TechArt, Unity, WebGL, Blender, Python.  
\item \textbf{Organización:} Mover repos irrelevantes (ej. repositorios personales básicos) a una organización diferente o marcarlos como “archived”.  
\item \textbf{Actividad social:} Contribuir a otros repos (como issues/respuestas en Unity o StackOverflow) y subir al menos un issue o PR público (demostrar colaboración).  
\item \textbf{Herramientas de análisis:} Usar GitHub Copilot o Advanced Security insights para mejorar la calidad del código y documentar cambios.  
\item \textbf{Cumplimiento legal:} Confirmar que no hay violación de licencias de terceros en los repos. Limpiar cualquier asset con copyright (usar propias texturas).  
\item \textbf{Metadata:} En profile de GitHub, incluir bio breve con rol deseado y ubicación.  
\end{itemize}




\subsection{15. Estrategia de portafolio}


\begin{itemize}
\item \textbf{Contenido actual:} Sugerencias: sitio web MVP, case studies (TwinSight, ARA, Retrato), uso de ArtStation, PDF one-pager, demo reel/video de TwinSight con guión y shots.  
\item \textbf{Faltantes / Correcciones:}  
\item \textbf{Sitio portfolio:} Definir tecnología (HTML/CSS/JS estático, React, o Hugo) y hosting (GitHub Pages, Netlify). Prototipo fácil de actualizar.  
\item \textbf{Estructura web:} Incluir páginas separadas: Home (resumen breve + enlace a LinkedIn/GitHub), Portafolio (videos/demos), Blog/Artículos (1-2 posts técnicos), Contacto.  
\item \textbf{Case studies completas:} Para cada proyecto principal, detallar: título, rol, tecnologías, desafíos, soluciones, resultados con métricas (ej. “Reducción a 95k polígonos – validado con 12 usuarios”). Incluir referencias a contexto académico (TwinSight es tesis UNAD).  
\item \textbf{AI News Aggregator:} Si se publica (ver punto 18), prever sección propia o mencionar en “Otros proyectos relevantes”.  
\item \textbf{FitApp-Free:} Evaluar su lugar: si se hace público, incluir como proyecto web completo; si no, mencionarlo brevemente en CV pero evitar en portafolio principal.  
\item \textbf{ArtStation:} Completar perfil: subir retrato Blender con breakdown (wireframe, texturas). Posiblemente crear “Technical breakdown” para cada artwork (topología, shaders, grooming). Incluir al menos 3-5 piezas (pueden ser variaciones del retrato).  
\item \textbf{PDF Portfolio:} Crear un PDF/slide deck breve (4-6 slides) que contenga el resumen de TwinSight y ARA, listo para enviar por email. Diseñarlo profesional (InDesign, Canva).  
\item \textbf{Demo reel/video:} Planificar grabación de la demo: escribir guión en inglés, preparar gráficos. Definir formato (Full HD, narración en voz o texto).  
\item \textbf{Organización de imágenes:} Crear repositorios de capturas/screenshots (alta calidad) para usar en CV, LinkedIn, portafolio. Etiquetarlos (TwinSight-assembly, ARA-architecture, etc).  
\item \textbf{Accesibilidad:} Considerar subtítulos o descripciones para el video/demo.  
\item \textbf{Mantenimiento:} Plan de actualizar el portafolio 2 veces al año con proyectos nuevos.  
\end{itemize}




\subsection{16. Plantilla de caso de estudio – TwinSight X500}


\begin{itemize}
\item \textbf{Faltantes / Requerido:} Crear un outline completo del caso de estudio:  
\item \textbf{Antecedentes:} Describir el problema (dificultad de entender manuales CAD de un dron). Contexto: \emph{“Proyecto de tesis UNAD 2026”}.  
\item \textbf{Objetivos:} Qué se quería lograr (mejorar comprensión, interactividad, pruebas de usabilidad).  
\item \textbf{Tecnologías:} Unity WebGL, C\#, URP, Blender (retopología), UI Toolkit, mediciones SUS/NASA-TLX.  
\item \textbf{Rol del candidato:} Detallar tareas propias (programación, UI, optimización, modelado de hélices, coordinación del estudio).  
\item \textbf{Features clave:} Lista de funcionalidades (exhibición ensamblaje, vistas explosivas, cross-section, modos visuales X-ray, blueprint, térmico).  
\item \textbf{Proceso:} Resumen de pipeline (de CAD/STEP a modelos optimizados, uso de Blender y Unity).  
\item \textbf{Resultados:} Métricas concretas (reducción de 6.5M a 95k tris, SUS 91.88 vs 59 de base, NASA-TLX 8.69 vs 19.89).  
\item \textbf{Impacto:} Qué demostró este proyecto (validación de usabilidad, potencial comercial).  
\item \textbf{Lecciones aprendidas:} Qué se mejoraría (p.ej. mayor optimización WebGL, integrar feedback más, thermal dinámico).  
\item \textbf{Medios:} Incluir diagramas o imágenes (bloque de Unity, esquemas de uso, foto del dron).  
\item \textbf{Palabras clave ATS:} Unity, WebGL, CAD, 3D Visualization, Technical Artist.  
\item \textbf{Evidencias:} Enlaces a demo WebGL si existe, screenshots, link repo GitHub.  
\end{itemize}




\subsection{17. Plantilla de caso de estudio – ARA Framework}


\begin{itemize}
\item \textbf{Faltantes / Requerido:} Desarrollar similar estructura:  
\item \textbf{Problema:} Investigación académica consume mucho tiempo (búsqueda bibliográfica, organización).  
\item \textbf{Solución:} Sistema multi-agente en Python (LangGraph, LangChain, Redis) para automatizar generación de reportes.  
\item \textbf{Rol propio:} Diseñó la arquitectura de agentes (nicho, investigación, síntesis), desarrolló la integración con APIs (Semantic Scholar, MarkItDown), coordinó prototipos y pruebas.  
\item \textbf{Tecnologías:} Python 3.12, LangGraph, Groq LLaMA 3.3, LangChain, Redis, Playwright, Semantic Scholar API.  
\item \textbf{Estado:} Mencionar que es prototipo (“proof-of-concept”), nivel de completitud actual (p.ej. agentes 1–3 funcionales, se planea mejorar).  
\item \textbf{Resultados iniciales:} Ejemplo de output (resúmenes de papers, nicho propuesto) para demostrar funcionalidad.  
\item \textbf{Desafíos:} Velocidad, precisión del LLM, volumen de datos.  
\item \textbf{Lecciones:} Necesidad de curación humana, escalabilidad de agentes, preguntas futuras (fine-tuning LLMs).  
\item \textbf{Imágenes:} Diagrama de arquitectura multi-agente, capturas del flujo de trabajo.  
\item \textbf{Visibilidad:} Link a código y breve demo (si se crea un ejemplo público).  
\end{itemize}




\subsection{18. Recomendación – AI News Aggregator}


\begin{itemize}
\item \textbf{Faltantes / Requerido:} Decidir si incluir o descartar:  
\item \textbf{Estado actual:} El proyecto está privado y requiere ajustes. Determinar si vale la pena completarlo (posible demostar habilidad full-stack con AI).  
\item \textbf{Viabilidad:} Evaluar tiempo necesario para dejarlo publicable (quizá un MVP con fuente de noticias pública y análisis básico).  
\item \textbf{Posicionamiento:} Si se publica, presentarlo como “AI-driven content platform” en CV y portafolio secundario. Resaltar tecnología Python/React/AI.  
\item \textbf{Completar readme/demo:} Si se hace público, escribir README con casos de uso (e.g. “agrega noticias de IA de múltiples fuentes y las analiza”).  
\item \textbf{Incluir en portfolio:} Solo como “otro proyecto”, no primer plano. Si no, archivar privado y no mencionarlo.  
\item \textbf{CV:} Incluir 1-2 bullets destacando habilidades tecnológicas (FastAPI, React, extracción automática).  
\end{itemize}




\subsection{19. Recomendación – FitApp-Free}


\begin{itemize}
\item \textbf{Faltantes / Requerido:} Evaluar relevancia:  
\item \textbf{Objetivo:} Es proyecto personal de backend/fitness. Determinar si respalda el enfoque técnico-art o distrae.  
\item \textbf{Si continuar:} Podría mostrar habilidades de app móvil y lógica de programación, útil para full-stack secundario. Subirlo parcialmente (screens/demo de UI).  
\item \textbf{Si no:} Mantener privado y no incluir en los principales materiales.  
\item \textbf{Menciones:} Si se completa, agregar bullet CV: “Desarrollé backend de app de fitness con seguimiento de progreso de entrenamiento en Python/Flask y PostgreSQL.”  
\item \textbf{Portafolio:} Solo incluir en sección “Otros proyectos de software” con breve descripción.  
\end{itemize}




\subsection{20. Recomendación – Retrato hiperrealista en Blender}


\begin{itemize}
\item \textbf{Contenido actual:} Se incluyen viñetas.  
\item \textbf{Faltantes / Correcciones:}  
\item \textbf{Mostrar en ArtStation:} Subir la imagen final y desglosar técnica: etapas de modelado, sculpt, texturizado, grooming, shading y luces.  
\item \textbf{Case study:} En portfolio, incluir sección de “Character Pipeline”: describir flujo de trabajo de topología, UVs, NODES usados en shading, post-procesado.  
\item \textbf{Conectar con Technical Art:} Explicar que demuestra “pipeline de personaje de alta calidad” (competencia en retopología, materiales complejos).  
\item \textbf{Capturas:} Incluir wireframes y mapas de normales/materiales en portafolio/CV (para destacar habilidad técnica).  
\item \textbf{CV:} Usar bullets ya escritos, asegurarse de traducir a inglés consistente.  
\end{itemize}




\subsection{21. Plan de demo reel/video (60–120s)}


\begin{itemize}
\item \textbf{Faltantes / Requerido:}  
\item \textbf{Guión detallado:} Escribir guión en inglés: introducción breve (texto/audio): objetivo del proyecto; destacar 2-3 demos clave.  
\item \textbf{Secuencia de escenas:} Listar shots con tiempos aproximados:  
\end{itemize}

\begin{enumerate}
\item \emph{Contexto}: Texto “Interactive 3D visualization for drone assembly” + vista general drone (5s).  
\item \emph{UI/Interactividad}: Mostrar selección de pieza y panel de info (5s).  
\item \emph{Explotar}: Exploded view del dron (5s).  
\item \emph{Clipping}: Section view (5s).  
\item \emph{Modos Visuales}: Animación cambiando a X-Ray/Blueprint/Termal (15s total).  
\item \emph{Mobile/WebGL}: Interfaz en navegador móvil (5s).  
\item \emph{Estadísticas/UI final}: Mostrar datos SUS/TLX o gráfico rápido (5s).  
\end{enumerate}

\begin{itemize}
\item \textbf{Narración/subtítulos:} Preparar frases clave: “Optimized drone CAD data for real-time Unity visualization…” (inglés claro).  
\item \textbf{Herramientas:} Plan de grabación (Unity Recorder para video, edición en Premiere o DaVinci).  
\item \textbf{Formato:} Exportar en 1080p, considerar versión de 60s vertical (mobile).  
\end{itemize}




\subsection{22. Lista de capturas/screenshots técnicos}


\begin{itemize}
\item \textbf{Faltantes / Requerido:}  
\item \textbf{TwinSight:} Capturas de: dron completo, dron wireframe, Unity editor mostrando jerarquía, ventana de inspector de pieza, modos (X-Ray, Blueprint), gráficos de datos (tabla reducciones).  
\item \textbf{Blender:} Imágenes de: modelo high-poly vs retopology, UV layout, wireframe sobre render, materiales.  
\item \textbf{ARA Framework:} Diagrama de agentes, captura de un output de texto sintetizado (si disponible).  
\item \textbf{Generales:} Logo del proyecto (si se hace) o un diagrama 3D representativo.  
\item \textbf{Formatos:} Guardar al menos 1920×1080, nombrar archivos descriptivos (e.g. TwinSight\_exploded.png).  
\item \textbf{Uso:} Incluir estas imágenes en el portfolio web y en documentos PDF.  
\end{itemize}




\subsection{23. Banco de palabras clave}


\begin{itemize}
\item \textbf{Faltantes / Requerido:}  
\item \textbf{Keywords inglés:} Unity, C\#, Real-Time 3D, WebGL, URP, Interactive Visualization, Technical Artist, GPU Optimization, Shader Graph, Blender, CAD, Digital Twin, XR/AR/VR, LLM, AI Tools, Python, FastAPI, .NET, IGRA.  
\item \textbf{Keywords español (para adaptaciones locales):} Artista Técnico, 3D Interactivo, Visualización Técnica, Optimización 3D, Programador Unity, Inteligencia Artificial, Automación Python.  
\item \textbf{Sinónimos y variantes:} E.g. en inglés: “3D Visualization Developer”, “VR Developer”, “Simulation Engineer”.  
\item \textbf{Palabras de impacto:} “Optimization”, “Real-Time Rendering”, “Low-Poly Workflow”, “Interactive UI”, “Data Visualization”, “Digital Twin”.  
\item \textbf{Aspectos blandos:} “Problem-solving”, “Agile development” (para sistemas corporativos).  
\item \textbf{Tecnologías secundarias:} “Git/GitHub”, “React”, “FastAPI”, “SQL”, “Blender Modeling”, “Shader Programming”.  
\item \textbf{Instrucciones:} Usar este banco para CV/LinkedIn/GitHub descripciones y títulos de repositorios para ser encontrado por reclutadores.  
\end{itemize}




\subsection{24. Lista de empresas/objetivos (≥100)}


\begin{itemize}
\item \textbf{Contenido actual:} Ejemplos puntuales por categoría (juegos, VR, defensa, etc.) en párrafos.  
\item \textbf{Faltantes / Correcciones:}  
\item \textbf{Ampliar lista:} Elaborar tabla con al menos 100 empresas, incluyendo:  
\item \textbf{Videojuegos (AAA/indie):} Ubisoft, Electronic Arts, Nintendo, CD Projekt, Insomniac, Valve, Wargaming, Larian Studios, Naughty Dog, etc.  
\item \textbf{XR/Simulación:} Unity Technologies, Epic Games, Oculus (Meta), Magic Leap, Varjo, Nvidia Omniverse, Mitsubishi Electric (simulación VR), Boeing XR, Airbus A\&D.  
\item \textbf{Digital Twins/Industria:} Siemens Digital Industries, Bosch, PTC (Creo/ThingWorx), ANSYS, Siemens Energy (gems), Dassault Systèmes, AutoDesk.  
\item \textbf{Drons/Aeroespacial:} DJI (China), Lockheed Martin (división drones), Boeing, SpaceX (grafismo VR intern), Holybro (asistente), General Dynamics.  
\item \textbf{MedTech/Educación:} Siemens Healthineers, Philips, GE Healthcare (ARX labs), Boston Scientific (VR training), univers. Harvard VR lab, Unity Reflect para arquitectura.  
\item \textbf{Energía/Defensa:} Raytheon, Northrop Grumman (Simulators), DARPA XR programs, Boeing, Leonardo (Italia defense XR).  
\item \textbf{Consultorías TI:} Accenture (immersive tech), Globant, ThoughtWorks, Endava, Softtek, PSL (Medellín), UruIT, TCS, Wipro (digital twin projects).  
\item \textbf{Outsourcing LATAM:} Grupo RutaN, Enola Labs (CO), Levvel, Belatrix (LATAM offices).  
\item \textbf{Tech Startups:} Magic Leap, Formlabs, Hugging Face (AI), Neuralink (neurovisualización), startups de digital twin o medtech en Europa.  
\item \textbf{Empresas remotas:} GitLab, Automattic, InVision (tech artists), empresas fully-remote de XR.  
\item \textbf{Columnas requeridas:} País/región, tipo (sector), roles buscados, \% remoto (basado en ofertas actuales), idioma requerido, sueldo estimado, link carrera, Contactos de Linkedin, estrategia (aplicación directa, referidos).  
\item \textbf{Priorización:} Marcar empresas top-20 high-fit (p.ej. Unity, Siemens) con mayor prioridad.  
\item \textbf{Fuentes:} Usar LinkedIn/Glassdoor/jsereros para confirmar presencia de roles Unity/3D en esas compañías.  
\item \textbf{Notas legales:} Indicar para cada si permiten contratos B2B (p.ej. consultorías sí, product companies menos).  
\item \textbf{Diversificar:} Incluir al menos 10 empresas por cada región objetivo (EEUU, UE, LATAM, Asia remota).  
\item \textbf{Reclutadores/contactos:} Anotar head hunters especializados (p.ej. Arteria en España, Robert Half Tech).  
\end{itemize}




\subsection{25. Bolsas de empleo y comunidades}


\begin{itemize}
\item \textbf{Contenido actual:} Listado de portales generales (LinkedIn, Indeed, WWR), sitios específicos (Unity Connect, Polycount, ArtStation, Gamedev, CGSociety), XR (XR Bootcamp jobs, Slack), digital twin (grupos LinkedIn, Slack), freelance (Upwork, Toptal, Malt), LATAM locales, comunidades Discord y LinkedIn.  
\item \textbf{Faltantes / Correcciones:}  
\item \textbf{Fuentes adicionales:} Añadir portales como Glassdoor, StackOverflow Jobs (para remote), AngelList (startups remotas).  
\item \textbf{Comunidades:} Listar Discords/Slack específicos: “Tech Art Slack”, “VR/AR Association”, “Digital Twin Forum”. Grupos de Facebook o Meetups (e.g. XR development en Berlín).  
\item \textbf{LATAM:} Incluir comunidades hispanas: “Game Developers Colombia”, “Desarrolladores Unity LatAm”, foros de Spanish Game Dev.  
\item \textbf{Canales de reclutamiento:} Referir newsletters y accounts de Twitter/LinkedIn para ofertas (p.ej. @RemoteOk, @XRJobs).  
\item \textbf{Freelance específico:} Añadir sitios enfocados a tech (Gun.io, Hired, Freelancer.com, Turing.com).  
\item \textbf{Palabra clave en portales:} Ejemplo de filtros recomendados (ej. “Unity AND remote AND Latin America” en Remotive).  
\item \textbf{Eventos y conferencias:} Mencionar participación en conferencias virtuales (GDC, SIGGRAPH virtual) y cómo obtener redes.  
\item \textbf{Frecuencia de seguimiento:} Sugerir revisión periódica de estas fuentes (diariamente, semanalmente).  
\end{itemize}




\subsection{26. Sistema de aplicación y métricas}


\begin{itemize}
\item \textbf{Contenido actual:} Propuesta de rutina semanal (20–25h), apps por semana (5–10), categorías A/B/C de vacantes, personalización \textasciitilde{}50\%, uso de tracker, seguimiento a reclutadores después de 10 días, criterios de rechazo (salario < USD1.5k, condiciones inestables).  
\item \textbf{Faltantes / Correcciones:}  
\item \textbf{Criterios A/B/C detallados:} Definir ejemplos concretos para cada categoría, basados en los requisitos vs perfil (por ej. A: Unity roles remotos con 3D; C: juegos AAA si no hay experiencia).  
\item \textbf{Herramientas de automatización:} Mencionar uso de herramientas (e.g. LinkedIn Easy Apply, mail merge en e-mails) sin perder personalización.  
\item \textbf{Plantillas:} Preparar plantillas básicas de correo de aplicación/carta de presentación adaptables (en inglés).  
\item \textbf{KPIs:} Especificar métricas de seguimiento: \#postulaciones, \#entrevistas, ratio de respuesta, \#contactos nuevos. Revisión semanal y ajuste de estrategia según tasa de éxito.  
\item \textbf{Diversificación:} Asegurar incluir aplicaciones internas (LinkedIn) y externas (portales especializados).  
\item \textbf{Freelance como alternativa:} Mantener presencia en plataformas freelance como Plan B si las aplicaciones no rinden los frutos esperados.  
\item \textbf{Calendario:} Sugerir días dedicados a cada tarea: e.g. L-Mi aplicaciones, J-V networking/cursos.  
\item \textbf{Criterios de seguimiento:} Ej. si no hay respuesta en 2 semanas, enviar mensaje de seguimiento.  
\item \textbf{Revisión de ofertas:} Incluir chequeo de elementos críticos (contratista vs empleado, timmings) antes de aceptar entrevista.  
\end{itemize}




\subsection{27. Plantilla de seguimiento (Notion/Google Sheets)}


\begin{itemize}
\item \textbf{Faltantes / Requerido:}  
\item \textbf{Columnas recomendadas:} Empresa, Puesto, Link oferta, Contacto reclutador, Fecha de aplicación, Estado (Aplicada, Entrevista, Oferta, Rechazada), Fecha de última acción, Acción siguiente programada, Fuente (LinkedIn/Portal), Tipo A/B/C, Observaciones.  
\item \textbf{Fórmulas/KPIs:} Incluir cálculos automáticos: total apps, porcentaje de respuesta, \% entrevistas.  
\item \textbf{Etiquetas/Colores:} Por prioridad A/B/C (semáforo) para visualizar foco.  
\item \textbf{Ejemplo inicial:} Proveer un par de filas ejemplo ya completas (p.ej. aplicación a Unity y a empresa de simulación).  
\item \textbf{Instrucciones de uso:} Breve guía en la misma hoja sobre cómo actualizar tras cada avance.  
\end{itemize}




\subsection{28. Rutina semanal de búsqueda}


\begin{itemize}
\item \textbf{Contenido actual:} Mencionado implícitamente en el sistema de aplicación.  
\item \textbf{Faltantes / Correcciones:}  
\item \textbf{Horario recomendado:} Proponer bloques diarios: mañanas (1-2h enviar apps), tardes (1h networking/LinkedIn), fines de semana (revisar resultados, cursos).  
\item \textbf{Bloques de aprendizaje:} Incluir tiempo fijo semanal para cursos o práctica (p.ej. 5-7h/semana).  
\item \textbf{Networking periódico:} Programa semanal para contactarse con X personas nuevas (p.ej. martes y jueves dedicar 30 min a LinkedIn).  
\item \textbf{Revisión y ajuste:} Una sesión semanal (p.ej. domingo) para analizar métricas, actualizar tracker y plan del siguiente bloque de 7 días.  
\item \textbf{Balance personal:} Incluir descansos y ejercicio (recomendado 30 min/día) para evitar burnout en periodo de búsqueda.  
\item \textbf{Flexibilidad:} Ajustar rutina según respuestas; estar listo para cambiar enfoque (ej. más tiempo a aplicaciones si falta respuesta).
\end{itemize}




\subsection{29. Plan tentativo 30/60/90 días}


\begin{itemize}
\item \textbf{Contenido actual:} Detallado en el informe.  
\item \textbf{Faltantes / Correcciones:}  
\item \textbf{Más métricas:} Definir objetivos específicos (e.g. “30 días: CV listo + 10 contactos LinkedIn; 60 días: 15 apps enviadas, 3 entrevistas técnicas”).  
\item \textbf{Responsabilidades diarias:} Asignar tareas diarias concretas (p.ej. lunes revisar portales, miércoles editar video, viernes networking).  
\item \textbf{Revisiones periódicas:} Programar reuniones de revisión propia cada 2 semanas para evaluar progresos.  
\item \textbf{Plan de contingencia:} Qué hacer si no se logran metas intermedias (p.ej. reevaluar CV o roles si tras 30 días no hay entrevistas).  
\item \textbf{Preparación de recursos:} Asegurar terminar tesis y obtener títulos/constancias necesarios antes de aplicar.  
\item \textbf{Actualización continua:} Incluir “mantenerse informado de noticias del sector” (leer blogs, papers) para entrevistas.  
\end{itemize}




\subsection{30. Scripts de acercamiento (Inglés/Español)}


\begin{itemize}
\item \textbf{Contenido actual:} Plantillas para varios escenarios (reclutador, lead técnico, fundador, director Unity, rechazo oferta, explicar falta experiencia, AI, etc.) tanto en EN como ES.  
\item \textbf{Faltantes / Correcciones:}  
\item \textbf{Cobertura total:} Incluir plantillas para los casos faltantes solicitados:  
\item \emph{Game studio recruiter (EN/ES)} (similar a tech art, pero indicando pasión por juegos).  
\item \emph{Aplicación especulativa (EN)} a empresa sin oferta específica (hablar de interés en posibles posiciones).  
\item \emph{Petición de recomendación/Referral (EN)} a contacto interno (“I see you work at X, I’d love if you could refer me for Unity roles…”).  
\item \emph{Profesor/mentor (ES/EN)} solicitando carta de recomendación (“Spanish: Profesor X, agradezco su orientación…”).  
\item \emph{Seguimiento después de la entrevista (EN)} (“Thank you for your time, reaffirm interés, contact details”).  
\item \emph{Disponibilidad remota desde Colombia (ES/EN)} (“Disponible full-time remoto desde Colombia, con equipo propio”).  
\item \textbf{Personalización:} Asegurar que cada script tenga variables [Nombre, Empresa] claras.  
\item \textbf{Claridad en salario y contrato:} Mantener plantillas sobre negociación (ej. “Would you clarify if this is full-time or contract?”).  
\item \textbf{Traducciones exactas:} Revisar que los mensajes en español sean formales y correctos gramaticalmente.  
\item \textbf{Tono profesional:} Ajustar scripts para que no suenen genéricos ni presumidos; ejemplos: agradecer oportunidad, mostrar empatía al responder una negativa.  
\item \textbf{Guía de uso:} Acompañar cada plantilla con nota breve de cuándo aplicarla.  
\end{itemize}




\subsection{31. Preparación para entrevistas}


\begin{itemize}
\item \textbf{Contenido actual:} Breve mención de posibles preguntas técnicas Unity/WebGL y comunes. Tácticas de uso de IA discutidas. Estrategia de negociación salarial esbozada. Checklist de oferta.  
\item \textbf{Faltantes / Correcciones:}  
\item \textbf{Preguntas técnicas específicas:} Listar preguntas probables (e.g. “Explicar ciclo de evento en Unity”, “Cómo optimizar un Shader Graph”, “¿Qué es LOD en modelos 3D?”). Preparar respuestas concisas.  
\item \textbf{Preguntas de comportamiento:} Practicar el famoso “Cuéntame sobre ti”, “Fortalezas y debilidades”, “¿Por qué esta empresa?”. Tener ejemplos con el perfil.  
\item \textbf{Presentación del portfolio:} Preparar versión demo/tour de TwinSight para mostrar pantalla. Ensayar narración breve del proyecto.  
\item \textbf{Exposición de proyectos:} Preparar “storytelling” de cada proyecto clave usando método STAR (Situación, Tarea, Acción, Resultado).  
\item \textbf{Idiomas:} Practicar entrevista en inglés, posiblemente con un amigo o tutor para evaluar fluidez.  
\item \textbf{Técnicas de entrevista virtual:} Verificar setup (cámara, fondo, micrófono) y conexión antes de cada entrevista online.  
\item \textbf{Preguntas al reclutador:} Preparar de antemano qué preguntar (“team composition?”, “roadmap del producto?”, “oportunidad de capacitación?”).  
\item \textbf{Conocimiento de la empresa:} Investigar 1-2 noticias recientes o productos de la empresa entrevistante para demostrar interés.  
\item \textbf{Código/ejercicio:} Anticipar si puede haber test de código (pocas chances en TechArt, pero tal vez tomar un desafío Unity simple).  
\end{itemize}




\subsection{32. Estrategia de negociación}


\begin{itemize}
\item \textbf{Contenido actual:} Menciona llevar rangos de referencia, considerar modalidad (W2/1099/EOR), beneficios adicionales (vacaciones, flexibilidad).  
\item \textbf{Faltantes / Correcciones:}  
\item \textbf{Rango objetivo:} Establecer meta salarial mínima (USD 3k) y justificativo (mercado remoto Senior en LATAM rondan X).  
\item \textbf{Costos adicionales:} Incluir cálculo de deducciones sociales en Colombia para ofertar un neto aceptable.  
\item \textbf{Beneficios no monetarios:} Mencionar negociar días de vacaciones, patrocinios de formación, equipo de trabajo provisto, cobertura de salud privada (muy relevante para contratistas).  
\item \textbf{Contratación internacional:} Tener en cuenta propuestas de EOR: negociar pagos mensuales en USD, uso de moneda estable.  
\item \textbf{Estrategias:} Ej. “price anchoring”: dar primero un rango en que 3k+ sea aceptable. Practicar respuestas tipo “Mi expectativa está alrededor de [rango calculado] en este mercado.”  
\item \textbf{Documentación:} Llevar comparativas salariales impresas (Glassdoor/PayScale) si se entrevista presencial, para sustentar cifras.  
\item \textbf{Títulos de cargos:} Asegurar que al nivel de oferta corresponde su rol (p.ej. que “Engineer” suene a Senior).  
\item \textbf{Plan B:} Si no alcanza el salario deseado, negociar otros beneficios (bono por objetivo, stock, entrenamiento).  
\end{itemize}




\subsection{33. Lista de señales de alerta (Red Flags)}


\begin{itemize}
\item \textbf{Faltantes / Requerido:}  
\item \textbf{Condiciones contractuales:} No aceptar cláusulas confusas, pago exclusivamente en retenciones (alejarse de “no hay contrato claro”).  
\item \textbf{Firma de NDA excesivo:} Cuidado con requisitos de confidencialidad que bloqueen publicar trabajos profesionales.  
\item \textbf{Plazos de pago dudosos:} Empresas que tardan en pagar o pagan en moneda local devaluada.  
\item \textbf{Alta rotación:} Compañías con reclamos de despidos masivos recientes.  
\item \textbf{Entrevistas vagas:} Entrevistadores sin claridad en el rol o cronograma (puede indicar improvisación).  
\item \textbf{Oferta demasiado baja:} Cualquier oferta < USD 1.5k (meta inicial) por rol técnico completo.  
\item \textbf{Comunicaciones extrañamente informales:} Reclutadores pidiendo entrevistas solo por WhatsApp o sin confirmar detalles.  
\item \textbf{Falta de recorrido legal:} Ofertas en efectivo sin facturación posible.  
\item \textbf{Expertise desconocida:} Empresas no verificables o con mala reputación online.  
\item \textbf{Exceso de tests:} Pedir tarea enorme no pagada antes de negociar salario.  
\item \textbf{Contratos poco claros:} Indicar incongruencias entre lo prometido (beneficios, horario) y lo escrito.  
\end{itemize}




\subsection{34. Estrategias de búsqueda (conservadora, equilibrada, agresiva)}


\begin{itemize}
\item \textbf{Contenido actual:} Breve descripción de tres enfoques y recomendación de combinar estrategias equilibradas.  
\item \textbf{Faltantes / Correcciones:}  
\item \textbf{Estrategia conservadora:} Aplicar principalmente a empresas “seguras” y ofertas estables (consultoras, empresas consolidadas) que ofrecen trabajo remoto estándar. Enfocar en roles con fuerte correspondencia al perfil. Prepararse para aceptar salario inicial modesto mientras gana experiencia.  
\item \textbf{Estrategia equilibrada:} Mezclar candidaturas a medianas/pequeñas startups internacionales (posibles sueldos más altos, pero riesgo medio), y consultoras/estudios de tecnología en LATAM (salario menor, estabilidad, facilidad legal). Mantener portfolio en inglés y atención al reclutamiento global.  
\item \textbf{Estrategia agresiva:} Postular a roles en empresas top-tier (Google AR/Unity, NVIDIA, Meta VR) aun sin experiencia directa; networking intensivo con headhunters; contratación como consultor independiente con tarifas más altas. Requiere high volume de outreach y preparación para rechazos.  
\item \textbf{Comparativa:} Tabla breve de pros/contras para cada (salario vs riesgo vs tiempo de contratación).  
\item \textbf{Recomendación final:} Explicar por qué se adopta un mix equilibrado: balance entre oportunidades asequibles y aspiraciones altas, asegurando avance en cada etapa.  
\item \textbf{Asignación de recursos:} Recomendaciones numéricas (ej. 50\% aplicaciones enfoque conservador, 30\% equilibrado, 20\% agresivo) con revisión trimestral.  
\end{itemize}




\subsection{35. Brechas de habilidades (ROI) – roadmap}


\begin{itemize}
\item \textbf{Contenido actual:} Lista priorizada de habilidades a desarrollar (Unity profiling, WebGL optimización, Shader/HLSL, Three.js, React, XR, C\# patterns, Python FastAPI, LangGraph, Blender, video/editing, escritura técnica, idiomas, Agile).  
\item \textbf{Faltantes / Correcciones:}  
\item \textbf{Categorizar por prioridad:} Asignar corto (<3 meses), mediano (3–12 meses) y largo plazo (>12m) a cada habilidad.  
\item \textbf{Integración con cursos:} Relacionar cada ítem con cursos concretos (e.g. “Unity WebGL – curso oficial Unity o tutorial en línea”, “Shader/HLSL – Coursera).  
\item \textbf{Medición de progreso:} Sugerir pequeños proyectos o pruebas para validar cada habilidad (p.ej. mejorar perfil de un demo, subir un shader propio a GitHub).  
\item \textbf{Idiomas específico:} Para alemán/portugués, fijar meta mínima (hablar con fluidez básica en 1 año, con apps semanales).  
\item \textbf{Herramientas complementarias:} Git avanzado (branching, CI/CD) y metodologías ágiles (uso básico de JIRA/Trello).  
\item \textbf{Soft skills:} Añadir comunicación en inglés fluida y trabajo en equipo virtual.  
\item \textbf{Fuentes:} Citar cursos recomendados para cada área (platzi, Coursera, Rebelway).  
\item \textbf{Planificación:} Crear plan de estudio semanal con fechas objetivo (ej. “Terminar Shader Graph avanzado en 4 semanas”).  
\end{itemize}




\subsection{36. Conclusiones y próximos pasos}


\begin{itemize}
\item \textbf{Contenido actual:} Resume brevemente el plan: posicionamiento como Unity TA, enfoque en Networking, formalidades contractuales, metas 3–6 meses.  
\item \textbf{Faltantes / Correcciones:}  
\item \textbf{Resumen ejecutivo:} Destacar 3 hallazgos clave: (1) “Perfil técnico sólido pero junior, necesita visibilidad y experiencia práctica”, (2) “Objetivo salarial de 3k–6k requiere inicio en roles remotos bien remunerados con proyección”, (3) “Estrategia clara en CV/LinkedIn + networking es vital”.  
\item \textbf{Metas específicas:} Reiterar objetivos concretos (e.g. “Conseguir al menos 3 ofertas de entrevista en 3 meses, alcanzar USD 3k en 6-12 meses”).  
\item \textbf{Iteración continua:} Enfatizar la necesidad de revisar el plan cada mes usando las métricas del tracker, y ajustar variantes de CV/objetivos en base a feedback.  
\item \textbf{Mentalidad:} Recordar persistencia y aprendizaje constante (en cursos, idioma, proyectos) mientras se busca empleo.  
\item \textbf{Agradecimientos/documentación:} Notar que todo el material (CV, scripts, plan) estará disponible como recursos reutilizables.  
\item \textbf{Plan a largo plazo:} Vincular con objetivo 12-24 meses (p. ej. “avanzar hacia roles senior o migración EU tras mejorar perfil”).  
\end{itemize}




\subsection{37. Formación continua (Cursos y especializaciones)}


\begin{itemize}
\item \textbf{Faltantes / Requerido:} Incluir una sección con:  
\item \textbf{Cursos completados:} Resumir brevemente los cursos relevantes hechos (sin certificado) listados en el Excel (CG Cookie, Animation, Modeling, Compositing, Videoviva, Web interactiva, Autodesk, etc.) para mostrar habilidades autodidactas.  
\item \textbf{Cursos recomendados a corto plazo (0–6 meses):} Ej. cursos en línea de Unity avanzado (URP optimización, GPU instancing), Blender avanzado (Hard Ops, iluminación técnica), React/Three.js (para complementar WebGL), Git y prácticas ágiles. Incluir los cursos de Rebelway accesibles: Houdini Pyro, Nuke avanzado, Katana.  
\item \textbf{Mediano plazo (6–12 meses):} Especializaciones o bootcamps: p.ej. certificación Unity, curso de XR Foundation (ARKit/ARCore), diplomado en Visualización 3D (si existe). Aprendizaje de idiomas técnicos (alemán B1/A2).  
\item \textbf{Largo plazo (>12 meses):} Opciones académicas: maestrías o posgrados en campos afines (p.ej. Maestría en Realidad Virtual, Ingeniería de medios interactivos). Mencionar universidades UE o becas posibles.  
\item \textbf{Idiomas:} Detallar meta: Inglés fluido (nivel C2) continuo, empezar portugués básico (Grupo Beta Archivos sefardí), alemán como segunda prioridad con aplicabilidad mediana.  
\item \textbf{Metodología:} Integrar cursos en la rutina semanal, usar proyectos reales para práctica (p.ej. implementar reel de efectos con Houdini).  
\item \textbf{Fuentes:} Referenciar plataformas: Udemy, Coursera, Domestika, Rebelway, CGCookie; mencionar uso de webinars y tutoriales gratuitos.
\end{itemize}



